Audio generation, also known as neural vocoding, aims to reconstruct high-resolution waveforms from compressed representations such as Mel-spectrograms and discrete audio tokens. It serves as a crucial component in various applications, including text-to-speech (TTS)~\citep{neural_tts,vits,zipvoice} and music synthesis~\citep{neural_audio_synth,adversarial_audio_synth,gansynth}, directly impacting the perceptual experience of end users.

Deep generative models have demonstrated remarkable progress in this task. 
Generative Adversarial Networks (GANs)~\citep{gan,wgan} are the dominant approach~\citep{melgan,hifigan,bigvgan,vocos}, leveraging sophisticated discriminators~\citep{melgan, hifigan,mrd,multi-subband-gan} to capture multi-grained audio details and enabling high-fidelity results with efficient one-step inference. However, adversarial optimization in GANs often leads to slow convergence and the risk of mode collapse~\citep{gan_collapse}. Another line of research explores the use of diffusion models~\citep{difussion_2015,ddpm,ddim} for audio generation~\citep{diffwave,Wavegrad,priorgrad,fregrad}, offering stable training and strong generative performance. 
Recent diffusion-based works that integrate multi-band waveform processing~\citep{mbd, rfwave} further boost diffusion on audio data, enabling them to generate high-quality audio.
% Recent diffusion-based works that integrate multi-band waveform processing have achieved GAN-competitive generation quality.
Nevertheless, the strong results of diffusion models come at the cost of multi-step sampling, resulting in high computational demands. 

In this work, we propose Flow2GAN, which follows a two-stage training strategy that combines the strengths of both GAN and diffusion paradigms. In the first stage, we train the model with a Flow Matching objective~\citep{fm1,fm2}, a form of diffusion model~\citep{kingma2023understanding}, to learn robust generative capabilities with stable training. In the second stage, we construct few-step generators from the trained model and utilize GAN fine-tuning for finer-grained generation through adversarial learning. 
To be specific, considering audio's distinctive properties, we improve Flow Matching for audio modeling by: 1) reformulating the training objective as direct endpoint prediction, circumventing the challenges of velocity estimation when involving empty regions; 2) incorporating spectral energy-based loss weighting to emphasize perceptually important low-energy regions. 
Building upon the adapted Flow Matching, the trained model is able to generate clearly improved audio with just 2 sampling steps (Figure~\ref{fig:fm_sample}). This provides a strong foundation for the subsequent GAN fine-tuning stage, which effectively guides the initialized few-step models toward further refined audio quality. Moreover, we design a multi-branch ConvNeXt-based~\citep{convnext} network structure that operates on multi-resolution Fourier coefficients, serving as a powerful backbone with enhanced modeling capabilities compared to the previous single-branch version. 
Our experiments show that Flow2GAN generates high-quality audio across both Mel-spectrogram and audio token conditioning, while providing highly favorable quality-efficiency trade-offs compared to state-of-the-art GAN-based and Flow Matching-based methods.
Ablation studies further validate the effectiveness of each proposed component. 



